% =====================================================================
%  minggu-11.tex  --  Slide Pertemuan Minggu ke-11
%  Pengantar Kriptografi (MPFA014) -- Prodi Matematika FMIPA UBB
%
%  Topik silabus : Algoritma RSA
%  Sumber materi : buku_ajar/chapters/10 (Bab 10)
%  Sub-CPMK      : 1.4
%
%  Kompilasi dari dalam folder ppt/:
%      pdflatex -interaction=nonstopmode minggu-11.tex   (dua kali)
% =====================================================================

\documentclass[aspectratio=169,11pt]{beamer}
% =====================================================================
%  beamer-preamble.tex  --  prasarana dan gaya slide kuliah
%
%  Mata kuliah : Pengantar Kriptografi (MPFA014)
%  Program     : Program Studi Matematika, FMIPA, Universitas Bale Bandung
%
%  Berkas ini TIDAK memuat \documentclass. Setiap berkas minggu memuatnya
%  sendiri lalu menarik berkas ini:
%
%      \documentclass[aspectratio=169,11pt]{beamer}
%      % =====================================================================
%  beamer-preamble.tex  --  prasarana dan gaya slide kuliah
%
%  Mata kuliah : Pengantar Kriptografi (MPFA014)
%  Program     : Program Studi Matematika, FMIPA, Universitas Bale Bandung
%
%  Berkas ini TIDAK memuat \documentclass. Setiap berkas minggu memuatnya
%  sendiri lalu menarik berkas ini:
%
%      \documentclass[aspectratio=169,11pt]{beamer}
%      % =====================================================================
%  beamer-preamble.tex  --  prasarana dan gaya slide kuliah
%
%  Mata kuliah : Pengantar Kriptografi (MPFA014)
%  Program     : Program Studi Matematika, FMIPA, Universitas Bale Bandung
%
%  Berkas ini TIDAK memuat \documentclass. Setiap berkas minggu memuatnya
%  sendiri lalu menarik berkas ini:
%
%      \documentclass[aspectratio=169,11pt]{beamer}
%      \input{beamer-preamble}
%      \title{...} \subtitle{...} \setWeekInfo{Minggu n}{Topik}
%
%  Berkas minggu hanya berisi \setWeekInfo dan lingkungan frame.
%  Kompilasi dijalankan DARI DALAM folder ppt/ (lihat compile-ppt.bat),
%  karena gambar buku dirujuk dengan jalur relatif ../buku_ajar/.
% =====================================================================

% ---------------------------------------------------------------- bahasa
\usepackage[utf8]{inputenc}
\usepackage[T1]{fontenc}
\usepackage[indonesian]{babel}
\usepackage{lmodern}
\usepackage{microtype}

% ------------------------------------------------------------- matematika
\usepackage{amsmath,amssymb,mathtools}

% ----------------------------------------------------------------- tabel
\usepackage{booktabs}
\usepackage{array}
\usepackage{tabularx}

% --------------------------------------------------------------- grafik
\usepackage{tikz}
\usetikzlibrary{arrows.meta,positioning,shapes.geometric,calc,fit,backgrounds}

% ---------------------------------------------------------- kode program
\usepackage{listings}
\lstdefinestyle{slidecode}{
    basicstyle=\ttfamily\scriptsize,
    keywordstyle=\color{ubbblue}\bfseries,
    commentstyle=\color{ubbgreen}\itshape,
    stringstyle=\color{ubbred},
    numbers=left,
    numberstyle=\tiny\color{gray},
    numbersep=6pt,
    frame=single,
    rulecolor=\color{ubbblue!40},
    backgroundcolor=\color{ubbblue!3},
    breaklines=true,
    showstringspaces=false,
    tabsize=4,
    columns=fullflexible,
    keepspaces=true,
}
\lstset{style=slidecode}

% =====================================================================
%  WARNA INSTITUSI
% =====================================================================
\definecolor{ubbblue}{RGB}{12, 62, 120}      % warna utama
\definecolor{ubbteal}{RGB}{0, 118, 110}      % aksen kedua
\definecolor{ubbgreen}{RGB}{22, 122, 74}     % untuk blok contoh
\definecolor{ubbred}{RGB}{168, 42, 42}       % untuk blok peringatan
\definecolor{ubbgray}{RGB}{70, 76, 84}       % teks sekunder

% =====================================================================
%  IDENTITAS MATA KULIAH
% =====================================================================
\newcommand{\CourseName}{Pengantar Kriptografi}
\newcommand{\CourseCode}{MPFA014}
\newcommand{\CourseSemester}{Semester Ganjil}
\newcommand{\CourseProgram}{Program Studi Matematika}
\newcommand{\CourseFaculty}{Fakultas Matematika dan Ilmu Pengetahuan Alam}
\newcommand{\CourseInstitution}{Universitas Bale Bandung}
\newcommand{\CourseLecturer}{Cecep Suwanda, S.Si., M.Kom.}

% Label minggu, disetel tiap berkas lewat \setWeekInfo
\newcommand{\WeekLabel}{}
\newcommand{\WeekTopic}{}
\newcommand{\setWeekInfo}[2]{%
    \renewcommand{\WeekLabel}{#1}%
    \renewcommand{\WeekTopic}{#2}%
}

% =====================================================================
%  TEMA: header, footer, warna blok
% =====================================================================
\usetheme{default}
\setbeamertemplate{navigation symbols}{}
\setbeamertemplate{blocks}[rounded][shadow=false]

\setbeamercolor{palette primary}{bg=ubbblue, fg=white}
\setbeamercolor{palette secondary}{bg=ubbgray, fg=white}
\setbeamercolor{palette tertiary}{bg=ubbteal, fg=white}

\setbeamercolor{structure}{fg=ubbblue}
\setbeamercolor{normal text}{fg=black!92}
\setbeamercolor{alerted text}{fg=ubbred}

\setbeamercolor{block title}{bg=ubbblue, fg=white}
\setbeamercolor{block body}{bg=ubbblue!6, fg=black!92}
\setbeamercolor{block title alerted}{bg=ubbred, fg=white}
\setbeamercolor{block body alerted}{bg=ubbred!6, fg=black!92}
\setbeamercolor{block title example}{bg=ubbgreen, fg=white}
\setbeamercolor{block body example}{bg=ubbgreen!6, fg=black!92}

\setbeamercolor{itemize item}{fg=ubbblue}
\setbeamercolor{itemize subitem}{fg=ubbteal}
\setbeamercolor{enumerate item}{fg=ubbblue}

\setbeamerfont{headline}{size=\scriptsize}
\setbeamerfont{footline}{size=\scriptsize}
\setbeamerfont{block title}{size=\normalsize}

% Badan blok rata kiri (bukan rata kanan-kiri): teks slide lebih mudah dibaca
% dan baris terakhir yang pendek tidak lagi memicu "Underfull \hbox".
\addtobeamertemplate{block body begin}{\raggedright}{}
\addtobeamertemplate{block alerted body begin}{\raggedright}{}
\addtobeamertemplate{block example body begin}{\raggedright}{}

% Header: nama mata kuliah dan kode MK
\setbeamertemplate{headline}{%
    \begin{beamercolorbox}[wd=\paperwidth,ht=2.6ex,dp=1.2ex,leftskip=1em,rightskip=1em]{palette primary}%
        \usebeamerfont{headline}\CourseName\ \textperiodcentered\ \CourseCode%
        \hfill \CourseProgram%
    \end{beamercolorbox}%
}

% Footer: minggu dan nomor halaman
\setbeamertemplate{footline}{%
    \begin{beamercolorbox}[wd=\paperwidth,ht=2.4ex,dp=1.2ex,leftskip=1em,rightskip=1em]{palette secondary}%
        \usebeamerfont{footline}\WeekLabel\ \textperiodcentered\ \WeekTopic%
        \hfill \insertframenumber\,/\,\inserttotalframenumber%
    \end{beamercolorbox}%
}

% Halaman judul: judul, subjudul, lalu blok identitas mata kuliah
\setbeamertemplate{title page}{%
    \begin{centering}
        \vskip1.0em
        {\usebeamerfont{title}\usebeamercolor[fg]{title}\inserttitle\par}
        \vskip0.3em
        {\usebeamerfont{subtitle}\usebeamercolor[fg]{subtitle}\insertsubtitle\par}
        \vskip1.2em
        \textcolor{ubbblue}{\rule{0.72\textwidth}{0.8pt}}
        \vskip1.2em
        {\small
            \textbf{\CourseName}\ (\CourseCode)\par
            \vskip0.4em
            \CourseProgram\par
            \CourseFaculty\par
            \CourseInstitution\par
            \vskip0.6em
            \CourseSemester\ \textperiodcentered\ \WeekLabel\par
            \vskip0.8em
            Dosen Pengampu: \CourseLecturer\par
        }
        \vskip1.0em
        \textcolor{ubbblue}{\rule{0.72\textwidth}{0.4pt}}
    \end{centering}
}

% =====================================================================
%  PERINTAH BANTU
% =====================================================================
\newcommand{\R}{\mathbb{R}}
\newcommand{\N}{\mathbb{N}}
\newcommand{\Z}{\mathbb{Z}}
\newcommand{\code}[1]{\texttt{#1}}

% Rujukan gambar dari buku: \bookfig[lebar]{NN}{namaberkas.png}
% Jalur relatif terhadap folder ppt/ (dikompilasi dari sana).
\newcommand{\bookfig}[3][0.80\linewidth]{%
    \includegraphics[width=#1]{../buku_ajar/chapters/#2/figures/#3}%
}

% Keterangan sumber gambar, sepadan dengan \sumber di buku
\newcommand{\figsumber}[1]{%
    \par\vspace{1pt}{\tiny\itshape Sumber: #1\par}%
}

% Penanda blok materi yang sering dipakai di slide
\newcommand{\slidejudul}[1]{{\large\bfseries\textcolor{ubbblue}{#1}}\par\vspace{2pt}}
\newcommand{\kunci}[1]{\textcolor{ubbred}{\textbf{#1}}}

%      \title{...} \subtitle{...} \setWeekInfo{Minggu n}{Topik}
%
%  Berkas minggu hanya berisi \setWeekInfo dan lingkungan frame.
%  Kompilasi dijalankan DARI DALAM folder ppt/ (lihat compile-ppt.bat),
%  karena gambar buku dirujuk dengan jalur relatif ../buku_ajar/.
% =====================================================================

% ---------------------------------------------------------------- bahasa
\usepackage[utf8]{inputenc}
\usepackage[T1]{fontenc}
\usepackage[indonesian]{babel}
\usepackage{lmodern}
\usepackage{microtype}

% ------------------------------------------------------------- matematika
\usepackage{amsmath,amssymb,mathtools}

% ----------------------------------------------------------------- tabel
\usepackage{booktabs}
\usepackage{array}
\usepackage{tabularx}

% --------------------------------------------------------------- grafik
\usepackage{tikz}
\usetikzlibrary{arrows.meta,positioning,shapes.geometric,calc,fit,backgrounds}

% ---------------------------------------------------------- kode program
\usepackage{listings}
\lstdefinestyle{slidecode}{
    basicstyle=\ttfamily\scriptsize,
    keywordstyle=\color{ubbblue}\bfseries,
    commentstyle=\color{ubbgreen}\itshape,
    stringstyle=\color{ubbred},
    numbers=left,
    numberstyle=\tiny\color{gray},
    numbersep=6pt,
    frame=single,
    rulecolor=\color{ubbblue!40},
    backgroundcolor=\color{ubbblue!3},
    breaklines=true,
    showstringspaces=false,
    tabsize=4,
    columns=fullflexible,
    keepspaces=true,
}
\lstset{style=slidecode}

% =====================================================================
%  WARNA INSTITUSI
% =====================================================================
\definecolor{ubbblue}{RGB}{12, 62, 120}      % warna utama
\definecolor{ubbteal}{RGB}{0, 118, 110}      % aksen kedua
\definecolor{ubbgreen}{RGB}{22, 122, 74}     % untuk blok contoh
\definecolor{ubbred}{RGB}{168, 42, 42}       % untuk blok peringatan
\definecolor{ubbgray}{RGB}{70, 76, 84}       % teks sekunder

% =====================================================================
%  IDENTITAS MATA KULIAH
% =====================================================================
\newcommand{\CourseName}{Pengantar Kriptografi}
\newcommand{\CourseCode}{MPFA014}
\newcommand{\CourseSemester}{Semester Ganjil}
\newcommand{\CourseProgram}{Program Studi Matematika}
\newcommand{\CourseFaculty}{Fakultas Matematika dan Ilmu Pengetahuan Alam}
\newcommand{\CourseInstitution}{Universitas Bale Bandung}
\newcommand{\CourseLecturer}{Cecep Suwanda, S.Si., M.Kom.}

% Label minggu, disetel tiap berkas lewat \setWeekInfo
\newcommand{\WeekLabel}{}
\newcommand{\WeekTopic}{}
\newcommand{\setWeekInfo}[2]{%
    \renewcommand{\WeekLabel}{#1}%
    \renewcommand{\WeekTopic}{#2}%
}

% =====================================================================
%  TEMA: header, footer, warna blok
% =====================================================================
\usetheme{default}
\setbeamertemplate{navigation symbols}{}
\setbeamertemplate{blocks}[rounded][shadow=false]

\setbeamercolor{palette primary}{bg=ubbblue, fg=white}
\setbeamercolor{palette secondary}{bg=ubbgray, fg=white}
\setbeamercolor{palette tertiary}{bg=ubbteal, fg=white}

\setbeamercolor{structure}{fg=ubbblue}
\setbeamercolor{normal text}{fg=black!92}
\setbeamercolor{alerted text}{fg=ubbred}

\setbeamercolor{block title}{bg=ubbblue, fg=white}
\setbeamercolor{block body}{bg=ubbblue!6, fg=black!92}
\setbeamercolor{block title alerted}{bg=ubbred, fg=white}
\setbeamercolor{block body alerted}{bg=ubbred!6, fg=black!92}
\setbeamercolor{block title example}{bg=ubbgreen, fg=white}
\setbeamercolor{block body example}{bg=ubbgreen!6, fg=black!92}

\setbeamercolor{itemize item}{fg=ubbblue}
\setbeamercolor{itemize subitem}{fg=ubbteal}
\setbeamercolor{enumerate item}{fg=ubbblue}

\setbeamerfont{headline}{size=\scriptsize}
\setbeamerfont{footline}{size=\scriptsize}
\setbeamerfont{block title}{size=\normalsize}

% Badan blok rata kiri (bukan rata kanan-kiri): teks slide lebih mudah dibaca
% dan baris terakhir yang pendek tidak lagi memicu "Underfull \hbox".
\addtobeamertemplate{block body begin}{\raggedright}{}
\addtobeamertemplate{block alerted body begin}{\raggedright}{}
\addtobeamertemplate{block example body begin}{\raggedright}{}

% Header: nama mata kuliah dan kode MK
\setbeamertemplate{headline}{%
    \begin{beamercolorbox}[wd=\paperwidth,ht=2.6ex,dp=1.2ex,leftskip=1em,rightskip=1em]{palette primary}%
        \usebeamerfont{headline}\CourseName\ \textperiodcentered\ \CourseCode%
        \hfill \CourseProgram%
    \end{beamercolorbox}%
}

% Footer: minggu dan nomor halaman
\setbeamertemplate{footline}{%
    \begin{beamercolorbox}[wd=\paperwidth,ht=2.4ex,dp=1.2ex,leftskip=1em,rightskip=1em]{palette secondary}%
        \usebeamerfont{footline}\WeekLabel\ \textperiodcentered\ \WeekTopic%
        \hfill \insertframenumber\,/\,\inserttotalframenumber%
    \end{beamercolorbox}%
}

% Halaman judul: judul, subjudul, lalu blok identitas mata kuliah
\setbeamertemplate{title page}{%
    \begin{centering}
        \vskip1.0em
        {\usebeamerfont{title}\usebeamercolor[fg]{title}\inserttitle\par}
        \vskip0.3em
        {\usebeamerfont{subtitle}\usebeamercolor[fg]{subtitle}\insertsubtitle\par}
        \vskip1.2em
        \textcolor{ubbblue}{\rule{0.72\textwidth}{0.8pt}}
        \vskip1.2em
        {\small
            \textbf{\CourseName}\ (\CourseCode)\par
            \vskip0.4em
            \CourseProgram\par
            \CourseFaculty\par
            \CourseInstitution\par
            \vskip0.6em
            \CourseSemester\ \textperiodcentered\ \WeekLabel\par
            \vskip0.8em
            Dosen Pengampu: \CourseLecturer\par
        }
        \vskip1.0em
        \textcolor{ubbblue}{\rule{0.72\textwidth}{0.4pt}}
    \end{centering}
}

% =====================================================================
%  PERINTAH BANTU
% =====================================================================
\newcommand{\R}{\mathbb{R}}
\newcommand{\N}{\mathbb{N}}
\newcommand{\Z}{\mathbb{Z}}
\newcommand{\code}[1]{\texttt{#1}}

% Rujukan gambar dari buku: \bookfig[lebar]{NN}{namaberkas.png}
% Jalur relatif terhadap folder ppt/ (dikompilasi dari sana).
\newcommand{\bookfig}[3][0.80\linewidth]{%
    \includegraphics[width=#1]{../buku_ajar/chapters/#2/figures/#3}%
}

% Keterangan sumber gambar, sepadan dengan \sumber di buku
\newcommand{\figsumber}[1]{%
    \par\vspace{1pt}{\tiny\itshape Sumber: #1\par}%
}

% Penanda blok materi yang sering dipakai di slide
\newcommand{\slidejudul}[1]{{\large\bfseries\textcolor{ubbblue}{#1}}\par\vspace{2pt}}
\newcommand{\kunci}[1]{\textcolor{ubbred}{\textbf{#1}}}

%      \title{...} \subtitle{...} \setWeekInfo{Minggu n}{Topik}
%
%  Berkas minggu hanya berisi \setWeekInfo dan lingkungan frame.
%  Kompilasi dijalankan DARI DALAM folder ppt/ (lihat compile-ppt.bat),
%  karena gambar buku dirujuk dengan jalur relatif ../buku_ajar/.
% =====================================================================

% ---------------------------------------------------------------- bahasa
\usepackage[utf8]{inputenc}
\usepackage[T1]{fontenc}
\usepackage[indonesian]{babel}
\usepackage{lmodern}
\usepackage{microtype}

% ------------------------------------------------------------- matematika
\usepackage{amsmath,amssymb,mathtools}

% ----------------------------------------------------------------- tabel
\usepackage{booktabs}
\usepackage{array}
\usepackage{tabularx}

% --------------------------------------------------------------- grafik
\usepackage{tikz}
\usetikzlibrary{arrows.meta,positioning,shapes.geometric,calc,fit,backgrounds}

% ---------------------------------------------------------- kode program
\usepackage{listings}
\lstdefinestyle{slidecode}{
    basicstyle=\ttfamily\scriptsize,
    keywordstyle=\color{ubbblue}\bfseries,
    commentstyle=\color{ubbgreen}\itshape,
    stringstyle=\color{ubbred},
    numbers=left,
    numberstyle=\tiny\color{gray},
    numbersep=6pt,
    frame=single,
    rulecolor=\color{ubbblue!40},
    backgroundcolor=\color{ubbblue!3},
    breaklines=true,
    showstringspaces=false,
    tabsize=4,
    columns=fullflexible,
    keepspaces=true,
}
\lstset{style=slidecode}

% =====================================================================
%  WARNA INSTITUSI
% =====================================================================
\definecolor{ubbblue}{RGB}{12, 62, 120}      % warna utama
\definecolor{ubbteal}{RGB}{0, 118, 110}      % aksen kedua
\definecolor{ubbgreen}{RGB}{22, 122, 74}     % untuk blok contoh
\definecolor{ubbred}{RGB}{168, 42, 42}       % untuk blok peringatan
\definecolor{ubbgray}{RGB}{70, 76, 84}       % teks sekunder

% =====================================================================
%  IDENTITAS MATA KULIAH
% =====================================================================
\newcommand{\CourseName}{Pengantar Kriptografi}
\newcommand{\CourseCode}{MPFA014}
\newcommand{\CourseSemester}{Semester Ganjil}
\newcommand{\CourseProgram}{Program Studi Matematika}
\newcommand{\CourseFaculty}{Fakultas Matematika dan Ilmu Pengetahuan Alam}
\newcommand{\CourseInstitution}{Universitas Bale Bandung}
\newcommand{\CourseLecturer}{Cecep Suwanda, S.Si., M.Kom.}

% Label minggu, disetel tiap berkas lewat \setWeekInfo
\newcommand{\WeekLabel}{}
\newcommand{\WeekTopic}{}
\newcommand{\setWeekInfo}[2]{%
    \renewcommand{\WeekLabel}{#1}%
    \renewcommand{\WeekTopic}{#2}%
}

% =====================================================================
%  TEMA: header, footer, warna blok
% =====================================================================
\usetheme{default}
\setbeamertemplate{navigation symbols}{}
\setbeamertemplate{blocks}[rounded][shadow=false]

\setbeamercolor{palette primary}{bg=ubbblue, fg=white}
\setbeamercolor{palette secondary}{bg=ubbgray, fg=white}
\setbeamercolor{palette tertiary}{bg=ubbteal, fg=white}

\setbeamercolor{structure}{fg=ubbblue}
\setbeamercolor{normal text}{fg=black!92}
\setbeamercolor{alerted text}{fg=ubbred}

\setbeamercolor{block title}{bg=ubbblue, fg=white}
\setbeamercolor{block body}{bg=ubbblue!6, fg=black!92}
\setbeamercolor{block title alerted}{bg=ubbred, fg=white}
\setbeamercolor{block body alerted}{bg=ubbred!6, fg=black!92}
\setbeamercolor{block title example}{bg=ubbgreen, fg=white}
\setbeamercolor{block body example}{bg=ubbgreen!6, fg=black!92}

\setbeamercolor{itemize item}{fg=ubbblue}
\setbeamercolor{itemize subitem}{fg=ubbteal}
\setbeamercolor{enumerate item}{fg=ubbblue}

\setbeamerfont{headline}{size=\scriptsize}
\setbeamerfont{footline}{size=\scriptsize}
\setbeamerfont{block title}{size=\normalsize}

% Badan blok rata kiri (bukan rata kanan-kiri): teks slide lebih mudah dibaca
% dan baris terakhir yang pendek tidak lagi memicu "Underfull \hbox".
\addtobeamertemplate{block body begin}{\raggedright}{}
\addtobeamertemplate{block alerted body begin}{\raggedright}{}
\addtobeamertemplate{block example body begin}{\raggedright}{}

% Header: nama mata kuliah dan kode MK
\setbeamertemplate{headline}{%
    \begin{beamercolorbox}[wd=\paperwidth,ht=2.6ex,dp=1.2ex,leftskip=1em,rightskip=1em]{palette primary}%
        \usebeamerfont{headline}\CourseName\ \textperiodcentered\ \CourseCode%
        \hfill \CourseProgram%
    \end{beamercolorbox}%
}

% Footer: minggu dan nomor halaman
\setbeamertemplate{footline}{%
    \begin{beamercolorbox}[wd=\paperwidth,ht=2.4ex,dp=1.2ex,leftskip=1em,rightskip=1em]{palette secondary}%
        \usebeamerfont{footline}\WeekLabel\ \textperiodcentered\ \WeekTopic%
        \hfill \insertframenumber\,/\,\inserttotalframenumber%
    \end{beamercolorbox}%
}

% Halaman judul: judul, subjudul, lalu blok identitas mata kuliah
\setbeamertemplate{title page}{%
    \begin{centering}
        \vskip1.0em
        {\usebeamerfont{title}\usebeamercolor[fg]{title}\inserttitle\par}
        \vskip0.3em
        {\usebeamerfont{subtitle}\usebeamercolor[fg]{subtitle}\insertsubtitle\par}
        \vskip1.2em
        \textcolor{ubbblue}{\rule{0.72\textwidth}{0.8pt}}
        \vskip1.2em
        {\small
            \textbf{\CourseName}\ (\CourseCode)\par
            \vskip0.4em
            \CourseProgram\par
            \CourseFaculty\par
            \CourseInstitution\par
            \vskip0.6em
            \CourseSemester\ \textperiodcentered\ \WeekLabel\par
            \vskip0.8em
            Dosen Pengampu: \CourseLecturer\par
        }
        \vskip1.0em
        \textcolor{ubbblue}{\rule{0.72\textwidth}{0.4pt}}
    \end{centering}
}

% =====================================================================
%  PERINTAH BANTU
% =====================================================================
\newcommand{\R}{\mathbb{R}}
\newcommand{\N}{\mathbb{N}}
\newcommand{\Z}{\mathbb{Z}}
\newcommand{\code}[1]{\texttt{#1}}

% Rujukan gambar dari buku: \bookfig[lebar]{NN}{namaberkas.png}
% Jalur relatif terhadap folder ppt/ (dikompilasi dari sana).
\newcommand{\bookfig}[3][0.80\linewidth]{%
    \includegraphics[width=#1]{../buku_ajar/chapters/#2/figures/#3}%
}

% Keterangan sumber gambar, sepadan dengan \sumber di buku
\newcommand{\figsumber}[1]{%
    \par\vspace{1pt}{\tiny\itshape Sumber: #1\par}%
}

% Penanda blok materi yang sering dipakai di slide
\newcommand{\slidejudul}[1]{{\large\bfseries\textcolor{ubbblue}{#1}}\par\vspace{2pt}}
\newcommand{\kunci}[1]{\textcolor{ubbred}{\textbf{#1}}}


\title{Algoritma RSA}
\subtitle{Satu modulus, dua kunci, dan keamanan yang bersandar pada sulitnya memfaktorkan}
\setWeekInfo{Minggu 11}{Algoritma RSA}

\begin{document}

% ----------------------------------------------------------- halaman judul
\begin{frame}[plain]
    \titlepage
\end{frame}

% ------------------------------------------------------------------ peta
\begin{frame}{Peta Materi Pertemuan Ini}
    \begin{columns}[T,onlytextwidth]
        \begin{column}{0.56\textwidth}
            \slidejudul{Lima bagian}
            \begin{enumerate}
                \item \textbf{Gagasan dan properti} --- peran kunci publik
                      dan kunci privat, serta lambang yang dipakai RSA.
                \item \textbf{Landasan matematis} --- Fungsi Totient Euler,
                      Teorema Euler, dan penurunan rumus.
                \item \textbf{Pembangkitan kunci} --- dari dua bilangan prima
                      sampai pasangan $(e,n)$ dan $(d,n)$.
                \item \textbf{Enkripsi dan dekripsi} --- rumus dan contoh
                      angka satu blok serta pesan multi-blok.
                \item \textbf{Keamanan dan kelemahan} --- pemfaktoran, panjang
                      modulus, kriptografi hibrida, \textit{padding}, dan
                      \textit{man-in-the-middle}.
            \end{enumerate}
        \end{column}
        \begin{column}{0.42\textwidth}
            \begin{block}{Referensi}
                \footnotesize
                Bab 10 buku ajar; Munir, \textit{IF4020 Kriptografi}, ITB,
                2026 (modul 13)
            \end{block}
            \begin{exampleblock}{Pertanyaan pemandu}
                \footnotesize
                Mengapa kunci yang boleh dibaca semua orang tetap tidak
                membocorkan kunci rahasianya?
            \end{exampleblock}
        \end{column}
    \end{columns}
\end{frame}

% ---------------------------------------------------- capaian pembelajaran
\begin{frame}{Capaian Pembelajaran}
    \begin{block}{Sub-CPMK yang dibebankan}
        \footnotesize
        \textbf{Sub-CPMK 1.4} --- Mengenal algoritma kriptografi kunci
        publik, yaitu RSA, Diffie--Hellman, ElGamal, knapsack, dan Elliptic
        Curve Cryptography (ECC), beserta persoalan matematis yang
        mendasarinya.
    \end{block}
    \vspace{0.1cm}
    \begin{exampleblock}{Setelah pertemuan ini, mahasiswa mampu}
        \footnotesize
        \begin{itemize}
            \item menjelaskan pembagian tugas kunci publik dan kunci privat
                  pada RSA;
            \item menjelaskan properti setiap lambang RSA dan status
                  kerahasiaannya;
            \item menjelaskan peran Fungsi Totient Euler dan Teorema Euler
                  dalam penurunan rumus RSA;
            \item membangkitkan pasangan kunci dari dua bilangan prima kecil;
            \item menghitung enkripsi dan dekripsi satu blok serta pesan
                  multi-blok;
            \item menjelaskan mengapa keamanan RSA bersandar pada sulitnya
                  pemfaktoran;
            \item mengidentifikasi kelemahan RSA dan risiko
                  \textit{man-in-the-middle}.
        \end{itemize}
    \end{exampleblock}
\end{frame}

% =====================================================================
\section{Gagasan dan Properti RSA}
% =====================================================================

\begin{frame}{Dua Kunci, Dua Tugas}
    \footnotesize
    \begin{columns}[T,onlytextwidth]
        \begin{column}{0.48\textwidth}
            \begin{block}{Kunci publik $(e,n)$}
                \begin{itemize}
                    \item Dipakai pengirim untuk mengenkripsi pesan.
                    \item Boleh diketahui siapa saja, termasuk penyadap.
                    \item Tidak memuat $d$, sehingga tidak dapat dipakai
                          membuka cipherteks.
                \end{itemize}
            \end{block}
        \end{column}
        \begin{column}{0.48\textwidth}
            \begin{block}{Kunci privat $(d,n)$}
                \begin{itemize}
                    \item Dipakai pemilik untuk mendekripsi cipherteks.
                    \item Harus disimpan rahasia dan tidak pernah dikirim.
                    \item Dibentuk dari $p$, $q$, dan $\varphi(n)$ yang juga
                          rahasia.
                \end{itemize}
            \end{block}
        \end{column}
    \end{columns}
    \vspace{0.15cm}
    \begin{alertblock}{Satu modulus untuk dua kunci}
        \footnotesize
        $n$ muncul pada kedua kunci karena seluruh perhitungan dilakukan
        modulo $n$. Bila $p$, $q$, dan $\varphi(n)$ bocor, siapa pun dapat
        menghitung $d$; karena itu ketiganya tidak boleh disebarkan.
    \end{alertblock}
\end{frame}

\begin{frame}{Penemu Algoritma RSA}
    \begin{columns}[T,onlytextwidth]
        \begin{column}{0.56\textwidth}
            \footnotesize
            \begin{itemize}
                \item Ron Rivest, Adi Shamir, dan Len Adleman merumuskan
                      algoritma ini di MIT pada tahun 1976.
                \item Nama RSA disusun dari huruf awal nama belakang
                      ketiganya.
                \item RSA adalah algoritma kunci publik yang paling dikenal
                      dan paling banyak dipakai.
                \item Yang disebarkan hanyalah kunci publik; kunci privat
                      tetap di tangan pemiliknya.
            \end{itemize}
        \end{column}
        \begin{column}{0.40\textwidth}
            \centering
            \bookfig[1.0\linewidth]{10}{penemu-rsa.png}
            \figsumber{Munir, \textit{IF4020 Kriptografi}, ITB.}
        \end{column}
    \end{columns}
    \vspace{0.1cm}
    \begin{exampleblock}{Pergeseran gagasan}
        \footnotesize
        Kriptografi simetri memakai satu kunci untuk mengunci dan membuka.
        RSA memisahkan kedua tugas itu ke dua kunci berbeda.
    \end{exampleblock}
\end{frame}

\begin{frame}{Lambang dan Status Kerahasiaan}
    \footnotesize
    \begin{tabularx}{\linewidth}{@{}c >{\raggedright\arraybackslash}X c@{}}
        \toprule
        \textbf{Lambang} & \textbf{Arti} & \textbf{Status} \\
        \midrule
        $p,\ q$      & Dua bilangan prima pembentuk modulus & Rahasia \\
        $n = p \cdot q$ & Modulus; panjangnya menentukan ukuran kunci & Umum \\
        $\varphi(n)$ & Fungsi Totient Euler, $(p-1)(q-1)$ & Rahasia \\
        $e$          & Eksponen enkripsi, disebut kunci publik & Umum \\
        $d$          & Eksponen dekripsi, disebut kunci privat & Rahasia \\
        $m$          & Plainteks yang dikirim & Rahasia \\
        $c$          & Cipherteks hasil enkripsi & Umum \\
        \bottomrule
    \end{tabularx}
    \vspace{0.15cm}
    \begin{block}{Dua pasangan yang beredar}
        \footnotesize
        Kunci publik adalah pasangan $(e,n)$; kunci privat adalah pasangan
        $(d,n)$. Hanya $n$ dan $e$ yang boleh disebarkan, sedangkan $p$, $q$,
        $\varphi(n)$, dan $d$ harus tetap rahasia.
    \end{block}
\end{frame}

\begin{frame}{Mengapa Modulus Dipakai Dua Kali}
    \footnotesize
    \begin{columns}[T,onlytextwidth]
        \begin{column}{0.52\textwidth}
            \begin{itemize}
                \item Enkripsi dan dekripsi sama-sama dihitung modulo $n$.
                \item Tanpa modulus yang sama, hasil dekripsi tidak kembali
                      kepada plainteks semula.
                \item Karena itu $n$ menjadi bagian dari kunci publik
                      sekaligus kunci privat.
                \item Yang membedakan kedua kunci hanyalah eksponennya:
                      $e$ untuk mengunci, $d$ untuk membuka.
            \end{itemize}
        \end{column}
        \begin{column}{0.44\textwidth}
            \begin{block}{Yang boleh dilihat penyadap}
                $n$ dan $e$ diketahui. Yang tidak diketahui adalah $p$, $q$,
                $\varphi(n)$, dan $d$. Selama $n$ sukar difaktorkan, $d$
                tetap terlindung.
            \end{block}
        \end{column}
    \end{columns}
    \vspace{0.15cm}
    \begin{alertblock}{Kesalahan yang sering terjadi}
        \footnotesize
        Menyebarkan $(e,n)$ adalah hal biasa; menyebarkan $p$, $q$, atau
        $\varphi(n)$ sama saja dengan menyerahkan kunci privat.
    \end{alertblock}
\end{frame}

% =====================================================================
\section{Landasan Matematis}
% =====================================================================

\begin{frame}{Fungsi Totient Euler}
    \footnotesize
    \begin{block}{Definisi}
        $\varphi(n)$ menghitung banyaknya bilangan antara $1$ dan $n$ yang
        relatif prima terhadap $n$, yaitu bilangan yang pembagi bersama
        terbesarnya dengan $n$ sama dengan satu.
    \end{block}
    \vspace{0.1cm}
    \begin{columns}[T,onlytextwidth]
        \begin{column}{0.50\textwidth}
            \begin{exampleblock}{Contoh dari buku: $\varphi(20) = 8$}
                \footnotesize
                Bilangan yang relatif prima terhadap $20$ adalah
                \[
                    1,\ 3,\ 7,\ 9,\ 11,\ 13,\ 17,\ 19 .
                \]
                Ada delapan bilangan, jadi $\varphi(20) = 8$.
            \end{exampleblock}
        \end{column}
        \begin{column}{0.46\textwidth}
            \begin{block}{Cara memeriksa}
                Setiap bilangan diuji dengan $\text{PBB}(a,20)$. Bilangan
                genap dan kelipatan lima gugur, sehingga yang tersisa tepat
                delapan bilangan.
            \end{block}
        \end{column}
    \end{columns}
\end{frame}

\begin{frame}{Totient untuk Modulus Dua Prima}
    \footnotesize
    \begin{block}{Rumus yang dipakai RSA}
        Bila $n = pq$ dengan $p$ dan $q$ prima, maka
        \[
            \varphi(n) = \varphi(p)\,\varphi(q) = (p-1)(q-1).
        \]
    \end{block}
    \begin{exampleblock}{Penerapan pada contoh buku}
        \footnotesize
        Untuk $p = 47$ dan $q = 71$:
        \[
            n = 47 \times 71 = 3337,
            \qquad
            \varphi(n) = 46 \times 70 = 3220 .
        \]
    \end{exampleblock}
    \begin{alertblock}{Mengapa $(p-1)(q-1)$}
        \footnotesize
        Di rentang $1$ sampai $pq$, yang tidak relatif prima hanyalah
        kelipatan $p$ atau $q$, sebanyak $q+p-1$. Sisanya
        $pq-(p+q-1)=(p-1)(q-1)$.
    \end{alertblock}
\end{frame}

\begin{frame}{Teorema Euler}
    \footnotesize
    \begin{block}{Pernyataan}
        Bila $a$ relatif prima terhadap $n$, yaitu $\text{PBB}(a,n) = 1$,
        maka
        \[
            a^{\varphi(n)} \equiv 1 \pmod{n}.
        \]
    \end{block}
    \vspace{0.1cm}
    \begin{columns}[T,onlytextwidth]
        \begin{column}{0.48\textwidth}
            \begin{alertblock}{Syarat tidak boleh dilupakan}
                Tanpa $\text{PBB}(a,n)=1$, kekongruenan itu tidak dijamin
                berlaku.
            \end{alertblock}
        \end{column}
        \begin{column}{0.48\textwidth}
            \begin{exampleblock}{Akibat untuk eksponen}
                Perpangkatan dapat direduksi modulo $\varphi(n)$: eksponen
                yang berbeda tetapi kongruen modulo $\varphi(n)$ memberi
                hasil yang sama.
            \end{exampleblock}
        \end{column}
    \end{columns}
    \vspace{0.1cm}
    \begin{block}{Tempatnya dalam RSA}
        Teorema ini yang membuat $m^{ed}$ dapat dikembalikan menjadi $m$.
        Seluruh kebenaran dekripsi RSA bertumpu padanya.
    \end{block}
\end{frame}

\begin{frame}{Menurunkan Rumus Enkripsi dan Dekripsi}
    \footnotesize
    Titik tolaknya Teorema Euler. Kedua ruas dipangkatkan $k$, lalu $a$
    diganti dengan plainteks $m$:
    \begin{align*}
        a^{\varphi(n)} &\equiv 1 \pmod{n}, \\
        a^{k\varphi(n)} &\equiv 1 \pmod{n}, \\
        m^{k\varphi(n)} &\equiv 1 \pmod{n}.
    \end{align*}
    Kedua ruas dikalikan $m$, sehingga
    \[
        m^{k\varphi(n)+1} \equiv m \pmod{n}.
    \]
    \begin{exampleblock}{Syarat pada $e$ dan $d$}
        \footnotesize
        Pilih $e$ dan $d$ sehingga $e \cdot d \equiv 1 \pmod{\varphi(n)}$,
        yang setara dengan $e \cdot d = k\varphi(n) + 1$.
    \end{exampleblock}
\end{frame}

\begin{frame}{Dari Penurunan ke Sepasang Rumus}
    \footnotesize
    \begin{block}{Bukti bahwa dekripsi membalik enkripsi}
        \[
            (m^{e})^{d} = m^{e \cdot d}
            = m^{\,k\varphi(n)+1} \equiv m \pmod{n}.
        \]
        Jadi enkripsi lalu dekripsi mengembalikan pesan semula.
    \end{block}
    \vspace{0.15cm}
    \begin{columns}[T,onlytextwidth]
        \begin{column}{0.48\textwidth}
            \begin{exampleblock}{Enkripsi}
                \[ c = m^{e} \bmod n \]
                Memakai kunci publik $(e,n)$.
            \end{exampleblock}
        \end{column}
        \begin{column}{0.48\textwidth}
            \begin{exampleblock}{Dekripsi}
                \[ m = c^{d} \bmod n \]
                Memakai kunci privat $(d,n)$.
            \end{exampleblock}
        \end{column}
    \end{columns}
    \vspace{0.1cm}
    \begin{alertblock}{Satu syarat tambahan}
        \footnotesize
        Pesan $m$ harus lebih kecil daripada $n$. Pesan yang lebih panjang
        dipecah menjadi blok sebelum dienkripsi.
    \end{alertblock}
\end{frame}

% =====================================================================
\section{Pembangkitan Pasangan Kunci}
% =====================================================================

\begin{frame}{Lima Langkah Pembangkitan Kunci}
    \footnotesize
    Pembangkitan kunci dijalankan satu kali oleh pemilik kunci.
    \begin{enumerate}
        \item Pilih dua bilangan prima $p$ dan $q$, sebaiknya $p \neq q$.
        \item Hitung modulus $n = p \cdot q$.
        \item Hitung $\varphi(n) = (p-1)(q-1)$.
        \item Pilih $e$ yang relatif prima terhadap $\varphi(n)$.
        \item Hitung $d$ dengan
              \[
                  e \cdot d \equiv 1 \pmod{\varphi(n)}
                  \quad \text{atau} \quad
                  d \equiv e^{-1} \pmod{\varphi(n)} .
              \]
    \end{enumerate}
    \begin{block}{Hasil akhir}
        \footnotesize
        Kunci publik $(e,n)$ dan kunci privat $(d,n)$. Nilai $p$, $q$, dan
        $\varphi(n)$ tidak lagi diperlukan setelah $d$ dihitung, tetapi
        tetap harus dimusnahkan atau disimpan rahasia.
    \end{block}
\end{frame}

\begin{frame}{Memilih $e$ dan Menghitung $d$}
    \footnotesize
    \begin{columns}[T,onlytextwidth]
        \begin{column}{0.48\textwidth}
            \begin{block}{Syarat $e$}
                \begin{itemize}
                    \item Relatif prima terhadap $\varphi(n)$.
                    \item Ganjil dan tidak terlalu kecil.
                    \item Nilai lazim $e = 65537$, dipakai pustaka RSA
                          karena aman dan efisien.
                \end{itemize}
            \end{block}
        \end{column}
        \begin{column}{0.48\textwidth}
            \begin{block}{Mencari $d$}
                Dari $e \cdot d = 1 + k\varphi(n)$ diperoleh
                \[
                    d = \frac{1 + k\,\varphi(n)}{e}.
                \]
                Nilai $k = 1, 2, 3, \dots$ dicoba sampai pembilangnya habis
                dibagi $e$.
            \end{block}
        \end{column}
    \end{columns}
    \vspace{0.1cm}
    \begin{exampleblock}{Mengapa $p \neq q$}
        \footnotesize
        Bila $p = q$, modulus berbentuk kuadrat sempurna sehingga $p$ dapat
        ditemukan lewat akar kuadrat biasa. Keamanan RSA hilang tanpa perlu
        pemfaktoran umum.
    \end{exampleblock}
\end{frame}

\begin{frame}{Contoh: Membangkitkan Kunci dari $p=47$ dan $q=71$}
    \footnotesize
    \begin{columns}[T,onlytextwidth]
        \begin{column}{0.52\textwidth}
            \begin{align*}
                n &= 47 \times 71 = 3337, \\
                \varphi(n) &= 46 \times 70 = 3220, \\
                \text{PBB}(79,3220) &= 1.
            \end{align*}
            Karena $79$ relatif prima terhadap $3220$, nilai $e = 79$ sah
            sebagai kunci publik. Kunci privat dicari dari
            \[
                d = \frac{1 + 3220k}{79}.
            \]
        \end{column}
        \begin{column}{0.46\textwidth}
            \begin{tabular}{@{}c c c c@{}}
                \toprule
                $k$ & $1+3220k$ & $/79$ & Bulat? \\
                \midrule
                $1$  & $3221$  & $40{,}77$ & tidak \\
                $2$  & $6441$  & $81{,}53$ & tidak \\
                $25$ & $80501$ & $1019$    & ya \\
                \bottomrule
            \end{tabular}
        \end{column}
    \end{columns}
    \vspace{0.1cm}
    \begin{exampleblock}{Pemeriksaan}
        \footnotesize
        $d = 80501/79 = 1019$, dan $79 \times 1019 = 80501 = 1 + 25 \times
        3220$. Jadi kunci publik $(79,3337)$ dan kunci privat $(1019,3337)$.
    \end{exampleblock}
\end{frame}

% =====================================================================
\section{Enkripsi dan Dekripsi}
% =====================================================================

\begin{frame}{Mengubah Pesan Menjadi Bilangan}
    \footnotesize
    \begin{block}{Pengodean}
        Setiap huruf diberi dua digit: $\text{A}=00$, $\text{B}=01$, sampai
        $\text{Z}=25$. Spasi tidak dikodekan dan ditambahkan kembali oleh
        penerima.
    \end{block}
    \begin{block}{Pemecahan menjadi blok}
        Deretan kode dipecah menjadi blok $m_1, m_2, \dots$ dengan syarat
        \[
            0 \le m_i \le n-1 .
        \]
        Setiap blok ditulis dengan panjang tetap, misalnya empat digit.
    \end{block}
    \begin{alertblock}{Nol di depan tidak boleh dibuang}
        \footnotesize
        Blok \code{0704} berbeda artinya dari \code{704} bila panjang blok
        tidak tetap. Menjaga nol di depan membuat batas antarblok tidak
        bergeser saat dirangkaikan kembali.
    \end{alertblock}
\end{frame}

\begin{frame}{Enkripsi dengan Kunci Publik}
    \footnotesize
    \begin{block}{Rumus per blok}
        Untuk setiap blok $m_i$ dan kunci publik $(e,n)$,
        \[
            c_i = m_i^{\,e} \bmod n .
        \]
    \end{block}
    \begin{itemize}
        \item Pengirim hanya memerlukan kunci publik penerima.
        \item Setiap blok dienkripsi secara terpisah dengan rumus yang sama.
        \item Cipherteks yang dikirim adalah deretan $c_1, c_2, \dots$ dengan
              panjang blok tetap.
        \item Ukuran blok harus lebih kecil daripada $n$ agar tidak ada
              informasi yang hilang.
    \end{itemize}
    \begin{exampleblock}{Sifat operasinya}
        \footnotesize
        Enkripsi adalah perpangkatan modular: mudah dihitung maju, tetapi
        akar pangkatnya sukar dicari tanpa kunci privat.
    \end{exampleblock}
\end{frame}

\begin{frame}{Dekripsi dengan Kunci Privat}
    \footnotesize
    \begin{block}{Rumus per blok}
        Pemilik kunci privat $(d,n)$ menghitung
        \[
            m_i = c_i^{\,d} \bmod n .
        \]
    \end{block}
    \begin{enumerate}
        \item Hitung kembali setiap blok plainteks dengan kunci privat.
        \item Rangkaikan seluruh $m_i$ sesuai urutan aslinya.
        \item Terjemahkan deretan angka itu menjadi huruf semula.
    \end{enumerate}
    \begin{alertblock}{Mengapa hanya pemilik kunci yang bisa}
        \footnotesize
        Tanpa $d$, cipherteks hanya deretan bilangan. Menghitung $d$
        menuntut $\varphi(n)$, dan menghitung $\varphi(n)$ menuntut
        pemfaktoran $n$.
    \end{alertblock}
\end{frame}

\begin{frame}{Alur Tiga Tahap RSA}
    \footnotesize
    \begin{center}
    \begin{tikzpicture}[
        kotak/.style={draw, rounded corners, align=center,
                      font=\scriptsize, inner sep=4pt},
        panah/.style={-{Latex[length=2.2mm]}, thick},
        kunci/.style={-{Latex[length=2.2mm]}, thick, dashed},
    ]
        \node[kotak, fill=ubbblue!8, text width=9.6cm] (gen) at (0,0)
            {Pembangkitan kunci: pilih $p,q$; hitung $n=pq$ dan
             $\varphi(n)$; pilih $e$; hitung $d$};
        \node[kotak, fill=ubbgreen!10, text width=2.8cm] (pub) at (-2.8,-1.4)
            {Kunci publik $(e,n)$};
        \node[kotak, fill=ubbred!8, text width=2.8cm] (priv) at (2.8,-1.4)
            {Kunci privat $(d,n)$};

        \node[kotak, fill=gray!8, text width=0.8cm] (m)   at (-5.9,-3.3) {$m$};
        \node[kotak, fill=orange!12, text width=2.6cm] (enc) at (-3.45,-3.3)
            {$c=m^{e}\bmod n$};
        \node[kotak, fill=gray!8, text width=0.7cm] (c)   at (-1.2,-3.3) {$c$};
        \node[kotak, fill=orange!12, text width=2.6cm] (dec) at (1.3,-3.3)
            {$m=c^{d}\bmod n$};
        \node[kotak, fill=gray!8, text width=0.8cm] (m2)  at (3.8,-3.3) {$m$};

        \draw[panah] (gen.south) -- ++(0,-0.35) -| (pub.north);
        \draw[panah] (gen.south) -- ++(0,-0.35) -| (priv.north);
        \draw[kunci] (pub.south) -- (enc.north);
        \draw[kunci] (priv.south) -- (dec.north);
        \draw[panah] (m) -- (enc);
        \draw[panah] (enc) -- (c);
        \draw[panah] (c) -- (dec);
        \draw[panah] (dec) -- (m2);
    \end{tikzpicture}
    \end{center}
    \vspace{0.05cm}
    \begin{block}{Bacaan gambar}
        \footnotesize
        Garis putus-putus adalah pemakaian kunci. Kunci publik mengunci,
        kunci privat membuka; keduanya bekerja pada modulus $n$ yang sama.
    \end{block}
\end{frame}

\begin{frame}{Contoh Inti: Enkripsi dan Dekripsi $m = 16$}
    \footnotesize
    Kunci publik Alice $(e,n) = (79,3337)$ dan kunci privatnya
    $(d,n) = (1019,3337)$.
    \vspace{0.1cm}
    \begin{enumerate}
        \item \textbf{Enkripsi oleh Bob:}
              \[
                  c = m^{e} \bmod n = 16^{79} \bmod 3337 = 1493 .
              \]
        \item \textbf{Pengiriman:} cipherteks $c = 1493$ dikirim melalui
              saluran terbuka.
        \item \textbf{Dekripsi oleh Alice:}
              \[
                  m = c^{d} \bmod n = 1493^{1019} \bmod 3337 = 16 .
              \]
        \item \textbf{Pemeriksaan:} plainteks pulih tepat $m = 16$.
    \end{enumerate}
    \begin{alertblock}{Dua hal yang mudah salah dibaca}
        \footnotesize
        Angka $16$ adalah bilangan bulat, bukan huruf. Cipherteks $1493$
        juga bukan pesan yang "disandikan huruf demi huruf"; seluruh blok
        diperlakukan sebagai satu bilangan.
    \end{alertblock}
\end{frame}

\begin{frame}{Contoh Multi-Blok: Mengodekan ``HELLO ALICE''}
    \footnotesize
    Aturan kode $\text{A}=00$ sampai $\text{Z}=25$; spasi diabaikan.
    \begin{center}
    \begin{tabular}{@{}l*{10}{c}@{}}
        \toprule
        Huruf & H & E & L & L & O & A & L & I & C & E \\
        \midrule
        Kode  & 07 & 04 & 11 & 11 & 14 & 00 & 11 & 08 & 02 & 04 \\
        \bottomrule
    \end{tabular}
    \end{center}
    Deretan kodenya menjadi
    \[
        m = 07041111140011080204 .
    \]
    Bilangan itu lebih besar daripada $n = 3337$, sehingga dipecah menjadi
    lima blok empat digit:
    \[
        m_1 = 0704,\quad m_2 = 1111,\quad m_3 = 1400,\quad
        m_4 = 1108,\quad m_5 = 0204 .
    \]
    \begin{exampleblock}{Pemeriksaan syarat blok}
        \footnotesize
        Kelima blok memenuhi $0 \le m_i \le 3336$, jadi masing-masing dapat
        dienkripsi tanpa pemotongan.
    \end{exampleblock}
\end{frame}

\begin{frame}{Enkripsi Lima Blok}
    \footnotesize
    Setiap blok dienkripsi dengan kunci publik Alice,
    $c_i = m_i^{79} \bmod 3337$.
    \vspace{0.1cm}
    \begin{center}
    \begin{tabular}{@{}c c c@{}}
        \toprule
        \textbf{Plainteks} & \textbf{Operasi} & \textbf{Cipherteks} \\
        \midrule
        $m_1 = \code{0704}$ & $704^{79} \bmod 3337$  & $c_1 = \code{0328}$ \\
        $m_2 = \code{1111}$ & $1111^{79} \bmod 3337$ & $c_2 = \code{0301}$ \\
        $m_3 = \code{1400}$ & $1400^{79} \bmod 3337$ & $c_3 = \code{2653}$ \\
        $m_4 = \code{1108}$ & $1108^{79} \bmod 3337$ & $c_4 = \code{2986}$ \\
        $m_5 = \code{0204}$ & $204^{79} \bmod 3337$  & $c_5 = \code{1164}$ \\
        \bottomrule
    \end{tabular}
    \end{center}
    \vspace{0.1cm}
    \begin{exampleblock}{Cipherteks yang dikirim}
        \footnotesize
        $C = 0328,\ 0301,\ 2653,\ 2986,\ 1164$. Setiap blok ditulis empat
        digit; nol di depan pada \code{0328} dan \code{0301} tetap
        dipertahankan.
    \end{exampleblock}
\end{frame}

\begin{frame}{Dekripsi Lima Blok}
    \footnotesize
    Alice memakai kunci privatnya, $m_i = c_i^{1019} \bmod 3337$.
    \vspace{0.1cm}
    \begin{center}
    \begin{tabular}{@{}c c c l@{}}
        \toprule
        \textbf{Cipherteks} & \textbf{Operasi} & \textbf{Pulih} & \textbf{Huruf} \\
        \midrule
        $c_1 = \code{0328}$ & $328^{1019} \bmod 3337$  & \code{0704} & H, E \\
        $c_2 = \code{0301}$ & $301^{1019} \bmod 3337$  & \code{1111} & L, L \\
        $c_3 = \code{2653}$ & $2653^{1019} \bmod 3337$ & \code{1400} & O, A \\
        $c_4 = \code{2986}$ & $2986^{1019} \bmod 3337$ & \code{1108} & L, I \\
        $c_5 = \code{1164}$ & $1164^{1019} \bmod 3337$ & \code{0204} & C, E \\
        \bottomrule
    \end{tabular}
    \end{center}
    \vspace{0.1cm}
    \begin{block}{Pesan kembali utuh}
        \footnotesize
        Deretan pulih $07041111140011080204$ diterjemahkan kembali menjadi
        ``HELLO ALICE''. Penerima menambahkan spasi yang tadi tidak
        dikodekan.
    \end{block}
\end{frame}

% =====================================================================
\section{Keamanan RSA}
% =====================================================================

\begin{frame}{Keamanan Bertumpu pada Pemfaktoran}
    \footnotesize
    \begin{block}{Satu langkah yang membuka semuanya}
        Bila $n$ berhasil difaktorkan menjadi $p$ dan $q$, maka
        \[
            \varphi(n) = (p-1)(q-1)
        \]
        dapat langsung dihitung.
    \end{block}
    \begin{block}{Kunci privat menyusul}
        Karena $e$ sudah diumumkan, nilai $d$ diperoleh dari
        \[
            e \cdot d \equiv 1 \pmod{\varphi(n)}
            \;\Longrightarrow\;
            d = \frac{1 + k\,\varphi(n)}{e}.
        \]
    \end{block}
    \begin{alertblock}{Kesimpulan}
        \footnotesize
        Memfaktorkan $n$ berarti memperoleh seluruh kunci privat. Inilah
        satu-satunya landasan keamanan RSA menurut bab ini.
    \end{alertblock}
\end{frame}

\begin{frame}{Biaya Pemfaktoran dan Panjang Modulus}
    \footnotesize
    \begin{block}{Kompleksitas algoritma tercepat saat ini}
        \[
            O\!\left(\exp\!\left(\sqrt[3]{\frac{64}{9}}\;
            b\,(\log b)^{2}\right)\right)
        \]
        untuk bilangan $n$ sepanjang $b$ bit. Sampai kini belum ada
        algoritma pemfaktoran bilangan besar dalam waktu polinomial;
        semakin panjang $n$, semakin lama pemfaktorannya.
    \end{block}
    \begin{exampleblock}{Perkiraan yang dikutip bab}
        \footnotesize
        \begin{itemize}
            \item $p$ dan $q$ lebih dari 100 digit: $n$ lebih dari 200
                  digit, diperkirakan perlu 4 miliar tahun.
            \item $n$ sepanjang 500 digit: diperkirakan perlu $10^{25}$
                  tahun.
            \item Perkiraan memakai komputer berkecepatan 1 milidetik per
                  operasi.
        \end{itemize}
    \end{exampleblock}
\end{frame}

\begin{frame}{Ukuran Modulus dan Rekor Pemfaktoran}
    \footnotesize
    \begin{columns}[T,onlytextwidth]
        \begin{column}{0.52\textwidth}
            \begin{tabular}{@{}l l@{}}
                \toprule
                \textbf{Panjang $n$} & \textbf{Keadaan menurut bab} \\
                \midrule
                $\le 256$ bit & beberapa jam dengan PC \\
                $\le 512$ bit & beberapa ratus komputer \\
                $> 1024$ bit & dianggap aman \\
                \bottomrule
            \end{tabular}
            \vspace{0.1cm}
            \begin{block}{Padanan desimal}
                $n$ 1024 bit setara 300 angka desimal; $p$ dan $q$
                masing-masing sekitar 154 angka desimal.
            \end{block}
        \end{column}
        \begin{column}{0.44\textwidth}
            \begin{exampleblock}{Sayembara 1977}
                \footnotesize
                Ketiga penemu RSA menantang pembaca \textit{Scientific
                American} memecahkan cipherteks RSA dengan hadiah \$100.
                Tantangan itu baru terpecahkan pada 1994, setelah 8 bulan
                kerja kolaborasi lewat internet.
            \end{exampleblock}
        \end{column}
    \end{columns}
    \vspace{0.1cm}
    \begin{alertblock}{Bacaan praktis}
        \footnotesize
        Rekor pemfaktoran terus bergerak. Modulus yang dulu dianggap besar
        dapat menjadi rentan hanya karena kemajuan perangkat keras.
    \end{alertblock}
\end{frame}

\begin{frame}{Ancaman Komputasi Kuantum}
    \footnotesize
    \begin{block}{Algoritma Shor}
        Komputer kuantum dapat memfaktorkan $n$ dalam waktu polinomial.
        Persoalan yang menjadi landasan RSA berubah dari "sukar" menjadi
        "dapat diselesaikan".
    \end{block}
    \vspace{0.1cm}
    \begin{itemize}
        \item Keamanan RSA bergantung pada pemfaktoran yang lambat.
        \item Memperpanjang modulus hanya menambah biaya, bukan mengubah
              sifat persoalan.
        \item Karena itu panjang kunci tidak dapat dipakai sebagai penangkal
              serangan kuantum.
    \end{itemize}
    \begin{alertblock}{Batas klaim}
        \footnotesize
        Ancaman ini bersifat prospektif: yang dibahas adalah akibat bila
        komputer kuantum berskala besar benar-benar terwujud. Selama itu
        belum terjadi, RSA masih dipakai dan masih aman terhadap penyerang
        klasik.
    \end{alertblock}
\end{frame}

% =====================================================================
\section{Kelemahan dan Man-in-the-Middle}
% =====================================================================

\begin{frame}{Kelemahan Pertama: RSA Lambat}
    \footnotesize
    \begin{columns}[T,onlytextwidth]
        \begin{column}{0.48\textwidth}
            \begin{alertblock}{Sebab}
                Perpangkatan modular dengan eksponen besar jauh lebih mahal
                daripada operasi blok pada cipher simetri. RSA lebih lambat
                daripada DES dan AES.
            \end{alertblock}
        \end{column}
        \begin{column}{0.48\textwidth}
            \begin{exampleblock}{Akibat praktis}
                RSA tidak dipakai untuk mengenkripsi berkas berukuran
                besar.
            \end{exampleblock}
        \end{column}
    \end{columns}
    \vspace{0.15cm}
    \begin{block}{Jalan keluar: kriptografi hibrida}
        \footnotesize
        \begin{enumerate}
            \item RSA mengenkripsi kunci sesi yang pendek dengan kunci
                  publik penerima.
            \item Pesan dienkripsi dengan cipher simetri, misalnya DES atau
                  AES, memakai kunci sesi itu.
            \item Penerima membuka kunci sesi dengan kunci privatnya, lalu
                  memakai kunci sesi untuk membuka pesan.
        \end{enumerate}
    \end{block}
\end{frame}

\begin{frame}{Kelemahan Kedua: Enkripsi Deterministik}
    \footnotesize
    \begin{alertblock}{Masalah bila pesan dienkripsi langsung}
        Tanpa \textit{padding}, pesan $m$ yang sama selalu menghasilkan
        cipherteks $c$ yang sama. Penyadap dapat mengenali pesan yang
        berulang, dan RSA terbuka terhadap serangan kriptografi.
    \end{alertblock}
    \begin{block}{OAEP sebagai penambal}
        OAEP (\textit{Optimal Asymmetric Encryption Padding}) diperkenalkan
        Bellare dan Rogaway pada 1994. Pesan diberi data tambahan sebelum
        dienkripsi:
        \[
            M \;\rightarrow\; \text{OAEP} \;\rightarrow\;
            \text{encoded message} \;\rightarrow\; \text{RSA}.
        \]
    \end{block}
    \begin{exampleblock}{Dua akibat yang diharapkan}
        \footnotesize
        Cipherteks menjadi tidak deterministik: enkripsi pesan yang sama dua
        kali memberi hasil berbeda. OAEP juga membuat RSA aman terhadap
        \textit{chosen-ciphertext attack}.
    \end{exampleblock}
\end{frame}

\begin{frame}{Peta Kelemahan RSA}
    \footnotesize
    \begin{tabularx}{\linewidth}{@{}>{\raggedright\arraybackslash}p{2.6cm}
                                    >{\raggedright\arraybackslash}X
                                    >{\raggedright\arraybackslash}X@{}}
        \toprule
        \textbf{Kelemahan} & \textbf{Pangkal} & \textbf{Akibat atau penanganan} \\
        \midrule
        Lambat & perpangkatan modular berukuran besar &
            tidak untuk data besar; dipakai hibrida \\
        \hline
        Deterministik & tanpa padding, $m$ sama memberi $c$ sama &
            pola terbaca; perlu OAEP \\
        \hline
        Modulus pendek & pemfaktoran makin murah &
            256 dan 512 bit tidak lagi memadai \\
        \hline
        Kunci tidak diautentikasi & kunci publik dapat diganti di jalan &
            \textit{man-in-the-middle} \\
        \hline
        Komputasi kuantum & algoritma Shor &
            modulus berapa pun tidak menolong \\
        \bottomrule
    \end{tabularx}
    \vspace{0.1cm}
    \begin{block}{Bacaan tabel}
        \footnotesize
        Hanya baris pertama dan kedua yang dapat ditambal di dalam algoritma
        RSA sendiri. Tiga baris sisanya menuntut keputusan di luar
        algoritma: panjang kunci, keaslian kunci, dan ancaman perangkat.
    \end{block}
\end{frame}

\begin{frame}{Serangan Man-in-the-Middle: Skenario}
    \footnotesize
    Alice dan Bob bertukar kunci publik melalui saluran terbuka. Carol
    mencegat keduanya.
    \vspace{0.1cm}
    \begin{enumerate}
        \item Alice mengirim kunci publiknya kepada Bob.
        \item Carol menahannya dan mengirim kunci publik Carol kepada Bob,
              seolah-olah itu kunci Alice.
        \item Carol melakukan hal yang sama terhadap Alice: kunci publik
              Carol dikirim seolah-olah itu kunci Bob.
        \item Alice dan Bob sama-sama percaya bahwa mereka memegang kunci
              lawan yang sah.
    \end{enumerate}
    \vspace{0.1cm}
    \begin{center}
    \begin{tikzpicture}[
        kotak/.style={draw, rounded corners, align=center,
                      font=\scriptsize, inner sep=4pt},
        panah/.style={-{Latex[length=2mm]}, thick},
    ]
        \node[kotak, fill=ubbgreen!10, text width=2.2cm] (A) at (0,0) {Alice};
        \node[kotak, fill=ubbred!10, text width=2.2cm] (C) at (3.6,0) {Carol};
        \node[kotak, fill=ubbgreen!10, text width=2.2cm] (B) at (7.2,0) {Bob};
        \draw[panah] (A) -- node[above, font=\tiny] {$K_A$} (C);
        \draw[panah] (C) -- node[above, font=\tiny] {$K_C$} (B);
        \draw[panah, dashed] (A.south) to[out=-35, in=215]
            node[below, font=\tiny] {Carol menyamar sebagai Bob} (B.south);
    \end{tikzpicture}
    \end{center}
\end{frame}

\begin{frame}{Man-in-the-Middle: Akibat dan Pelajarannya}
    \footnotesize
    \begin{block}{Yang dilakukan Carol}
        \begin{enumerate}
            \item Bob mengenkripsi pesan dengan kunci yang disangkanya kunci
                  Alice, padahal kunci Carol.
            \item Carol mendekripsi pesan itu dengan kunci privatnya dan
                  menyimpan salinannya.
            \item Carol mengenkripsi ulang pesan dengan kunci publik Alice
                  yang asli, lalu meneruskannya.
        \end{enumerate}
    \end{block}
    \begin{alertblock}{Mengapa tidak terdeteksi}
        \footnotesize
        Alice menerima pesan yang memang dapat dibuka dengan kunci
        privatnya, dan Bob yakin telah mengirim kepada Alice. Keduanya tidak
        melihat tanda apa pun bahwa Carol berada di tengah.
    \end{alertblock}
    \begin{exampleblock}{Pelajaran}
        \footnotesize
        RSA memberi kerahasiaan, bukan identitas. Keaslian kunci publik
        harus dipastikan melalui saluran lain yang dipercaya sebelum kunci
        itu dipakai.
    \end{exampleblock}
\end{frame}

% =====================================================================
\section{Rangkuman}
% =====================================================================

\begin{frame}{Rangkuman (1/2): Gagasan dan Prosedur}
    \footnotesize
    \begin{itemize}
        \item RSA memakai sepasang kunci: $(e,n)$ untuk enkripsi dan
              $(d,n)$ untuk dekripsi. Hanya $n$ dan $e$ yang disebarkan.
        \item $p$ dan $q$ prima rahasia, $n = pq$ umum, dan
              $\varphi(n) = (p-1)(q-1)$ rahasia.
        \item Teorema Euler, $a^{\varphi(n)} \equiv 1 \pmod n$ untuk
              $\text{PBB}(a,n)=1$, menopang penurunan rumus RSA.
        \item Dari $e \cdot d = k\varphi(n)+1$ diperoleh
              $m^{ed} \equiv m \pmod n$, sehingga dekripsi membalik
              enkripsi.
        \item Pembangkitan kunci: pilih $p,q$; hitung $n$ dan
              $\varphi(n)$; pilih $e$ yang relatif prima; hitung
              $d \equiv e^{-1} \pmod{\varphi(n)}$.
        \item Pesan dipecah menjadi blok $0 \le m_i \le n-1$, ditulis
              dengan panjang tetap agar nol di depan tidak hilang.
    \end{itemize}
\end{frame}

\begin{frame}{Rangkuman (2/2): Contoh, Keamanan, dan Kelemahan}
    \footnotesize
    \begin{itemize}
        \item Contoh buku: $p=47$, $q=71$ memberi $n=3337$,
              $\varphi(n)=3220$, $e=79$, dan $d=1019$.
        \item Enkripsi $m=16$ memberi $c = 16^{79} \bmod 3337 = 1493$;
              dekripsi $1493^{1019} \bmod 3337$ memulihkan $16$.
        \item ``HELLO ALICE'' dikodekan menjadi
              $07041111140011080204$, dipecah lima blok, dan dienkripsi
              menjadi $0328, 0301, 2653, 2986, 1164$.
        \item Keamanan bersandar pada pemfaktoran $n$: sekali $p$ dan $q$
              diketahui, $\varphi(n)$ dan $d$ menyusul.
        \item Kelemahan utamanya: lambat, deterministik tanpa padding,
              rentan bila modulus pendek, dan tidak menjamin keaslian kunci.
        \item \textit{Man-in-the-middle} berhasil karena kunci publik dapat
              diganti di jalan; penangkalnya adalah kepastian identitas
              pemilik kunci.
    \end{itemize}
\end{frame}

\begin{frame}{Latihan}
    \footnotesize
    \begin{enumerate}
        \item Dengan $p = 47$ dan $q = 71$, hitunglah $\varphi(n)$. Periksa
              bahwa $e = 17$ memenuhi syarat, lalu tentukan $d$.
        \item Jelaskan mengapa $p$, $q$, dan $\varphi(n)$ harus dirahasiakan,
              sedangkan $n$ boleh diketahui umum.
        \item Dengan kunci publik $(79,3337)$ dan kunci privat
              $(1019,3337)$, enkripsikan $m = 16$ lalu dekripsikan hasilnya.
        \item Mengapa pesan yang lebih panjang daripada modulus harus
              dipecah menjadi blok dengan syarat $0 \le m_i \le n-1$?
        \item Jelaskan mengapa OAEP membuat cipherteks RSA menjadi tidak
              deterministik.
        \item Uraikan jalannya \textit{man-in-the-middle} ketika Alice dan
              Bob bertukar kunci publik, lalu sebutkan apa yang diperlukan
              untuk mencegahnya.
    \end{enumerate}
\end{frame}

\begin{frame}{Penutup}
    \begin{block}{Yang harus diingat}
        \footnotesize
        \begin{itemize}
            \item RSA memisahkan tugas mengunci dan membuka ke dua kunci
                  berbeda.
            \item Kebenaran dekripsi berasal dari Teorema Euler dan syarat
                  $ed \equiv 1 \pmod{\varphi(n)}$.
            \item Keamanannya hanya sekuat kesulitan memfaktorkan $n$.
            \item RSA lambat, sehingga dipakai secara hibrida untuk
                  membungkus kunci sesi.
            \item Tanpa kepastian keaslian kunci publik, RSA tidak menahan
                  \textit{man-in-the-middle}.
        \end{itemize}
    \end{block}
    \vspace{0.15cm}
    \begin{exampleblock}{Pertemuan berikutnya}
        \textbf{Minggu 12 --- Diffie--Hellman dan ElGamal}: logaritma
        diskret, pertukaran kunci Diffie--Hellman, serangan pada DH, dan
        algoritma ElGamal.
    \end{exampleblock}
\end{frame}

\end{document}
